\documentclass[11pt]{article}
\pagestyle{plain}
\usepackage[margin=1in]{geometry}
\usepackage{amsmath,amssymb,amsthm,booktabs}
\usepackage[T1]{fontenc}
\usepackage[utf8]{inputenc}
\usepackage{lmodern}
\usepackage[hyphens]{url}
\usepackage[colorlinks=true,linkcolor=blue,citecolor=blue,urlcolor=blue]{hyperref}
\newtheorem{theorem}{Theorem}
\newtheorem{lemma}[theorem]{Lemma}
\newcommand{\scplx}{\operatorname{sc}}
\newcommand{\codim}{\operatorname{codim}}
\setlength{\emergencystretch}{2em}
\raggedbottom
\title{A monotone counterexample to a question on distributional complexity}
\author{P. Graham\\ \small\url{https://orcid.org/0009-0007-8493-1505}}
\date{October 24, 2026}
\hypersetup{pdfauthor={P. Graham},pdftitle={A monotone counterexample to a question on distributional complexity}}
\begin{document}
\maketitle
\begingroup
\renewcommand{\thefootnote}{}
\footnotetext{The proofs, the Lean formalization, the search code, the certificate checkers and the text were produced by AI models (from OpenAI and Anthropic), working under the author's direction.}
\endgroup
\begin{abstract}
For a Boolean function under a product measure, distributional algorithmic
complexity is the minimum expected number of adaptive queries needed to
determine its value. Distributional subcube partition complexity is the
minimum expected number of fixed coordinates in a monochromatic subcube
partition. K\"ohler-Schindler and Steif asked in 2024 whether these quantities
are equal for every monotone Boolean function. We answer their Question~2.1
in the negative with the explicit function
$f(x_0,\ldots,x_4)=x_1x_2x_3\lor x_0x_4(x_1\lor x_2\lor x_3)$.
For independent fair input bits its algorithmic complexity is $13/4$,
while its subcube partition complexity is $51/16$.
A finite recurrence proves the algorithmic lower bound, and an explicit
partition gives the strict inequality. An exact linear-programming dual
certificate proves the partition value is optimal. The strict inequality
is formalized in Lean; independent exact checkers verify the additional
dual certificate.
\end{abstract}

\section{Introduction}
A deterministic adaptive algorithm evaluates a Boolean function by asking for input bits,
with each query depending on previous answers. A subcube partition instead
assigns each input to a unique partial assignment that determines the
function. Every decision tree gives such a partition through its leaves,
but not every partition comes from a decision tree. We compare their
average costs when the bits are independent fair coins.

Following Laurin K\"ohler-Schindler and Jeffrey E.\ Steif~\cite{KS}, write
$\pi_p$ for the product measure with independent Bernoulli($p$) coordinates.
For $g:\{0,1\}^n\to\{0,1\}$, let $a(g,\pi_p)$ be the minimum expected
number of adaptive queries until $g$ is determined. A subcube fixes some
coordinates and leaves the others free; its codimension is the number
fixed. A partition is monochromatic for $g$ if $g$ is constant on each
cell. Write $\scplx(g,\pi_p)$ for the minimum expected codimension of the
cell containing the random input, over all such partitions.
The decision-tree partitions imply $a(g,\pi_p)\geq\scplx(g,\pi_p)$.

K\"ohler-Schindler and Steif's paper was first posted on August~23, 2024~\cite{KS}.
We quote their Question~2.1 from arXiv version~1:
\begin{quote}
Does $a(f,\pi_p)=sc(f,\pi_p)$ hold for every monotone Boolean function $f$
and for every $p\in(0,1)$?
\end{quote}
Here monotone means nondecreasing in each input coordinate. The following
example answers this question.

\begin{theorem}\label{thm:main}
The Boolean function
\begin{equation}\label{eq:f}
 f(x_0,x_1,x_2,x_3,x_4)
 =x_1x_2x_3\lor x_0x_4(x_1\lor x_2\lor x_3)
\end{equation}
is monotone and satisfies
\[
 a(f,\pi_{1/2})=\frac{13}{4}>\frac{51}{16}
 =\scplx(f,\pi_{1/2}).
\]
Products in~\eqref{eq:f} mean Boolean conjunction, and $\lor$ means disjunction.
\end{theorem}

The contribution is this small monotone example. K\"ohler-Schindler and
Steif already gave strict distributional gaps for the nonmonotone
ALL-EQUAL function on three inputs and the four-input function $h$
\cite[arXiv version~1, Example~2.4 and Section~2.3]{KS}.
The additional task is to preserve
monotonicity while making a partition cheaper than every adaptive query
strategy. For~\eqref{eq:f}, the main lower-bound obligation is to rule out
all first queries and their possible continuations, rather than compare
only a few natural algorithms. The finite recurrence below does so.
We make no claim of minimality in the number of inputs, of a gap at other
values of $p$, or of an asymptotic separation.

Earlier separations between decision trees and subcube partitions concern
other costs. Kothari, Racicot-Desloges and Santha~\cite{KRS} established an
asymptotic separation of randomized decision-tree complexity from
deterministic subcube partition complexity using iterated quaternary
majority. Those are worst-input costs, not the product-measure expectations
here. Arunachalam, Chakraborty, Kouck\'y, Saurabh and de Wolf
studied average unambiguous certificate complexity
\cite[version~2, Section~2, unambiguous certificate complexity]{ACKSW};
after relabeling
$\{-1,1\}$ as $\{0,1\}$, their quantity is exactly
$\scplx(g,\pi_{1/2})$. Our definitions and question come from
K\"ohler-Schindler and Steif, and these earlier works provide the surrounding
comparison framework.

\subsection*{Proof overview}
The formula for $f$ makes monotonicity immediate. Section~\ref{sec:alg}
evaluates the optimal adaptive cost by conditioning on the next query.
Section~\ref{sec:partition} gives a monochromatic partition of smaller
average codimension; this already answers Question~2.1. It then proves
that the partition is optimal by assigning a signed weight to each input.
The sum of the weights in every monochromatic subcube is at most that
subcube's total codimension cost, so summing over any partition gives the
required lower bound. Section~\ref{sec:checks} explains the finite checks
and their precise formalization boundary.

This work has not yet been reviewed by experts.

\section{The adaptive cost}\label{sec:alg}
We determine the best adaptive algorithm by a recurrence that considers
every available next query. A restriction $r\in\{0,1,*\}^5$ specifies
some coordinates and leaves the starred coordinates free. For example,
$00***$ fixes $x_0=x_1=0$. Let $D(r)$ be the minimum expected number of
additional queries for $f$, conditional on the input belonging to $r$,
under $\pi_{1/2}$. All unqueried coordinates remain independent fair bits.
An irrelevant coordinate can be omitted: its answer is independent of the
remaining coordinates, and choosing the cheaper of the two continuation
strategies without asking that query cannot increase the expected cost.
Let $r^{i,b}$ be the restriction obtained by replacing a free coordinate
$i$ with $b\in\{0,1\}$.

\begin{lemma}\label{lem:recurrence}
If $f$ is constant on $r$, then $D(r)=0$. Otherwise,
\begin{equation}\label{eq:recurrence}
 D(r)=1+\min_{i:r_i=*}
       \frac{D(r^{i,0})+D(r^{i,1})}{2}.
\end{equation}
\end{lemma}
\begin{proof}
A query to an already fixed coordinate gives no information and can be
omitted. If the value is not determined, the first remaining query must
therefore be to a free coordinate $i$. Each answer has conditional
probability $1/2$. The two continuations have expected costs at least
$D(r^{i,0})$ and $D(r^{i,1})$. Conversely, choosing a minimizing $i$ and
optimal continuations attains the right side. Induction on the number of
free coordinates proves the assertion and existence of the minima.
\end{proof}

The recurrence starts at fully specified inputs and works upward. Its
values after the first query are shown in Table~\ref{tab:branches}.
The two kinds of row reflect the symmetries exchanging $x_0$ with $x_4$
and permuting $x_1,x_2,x_3$. Here $D(x_i=b)$ denotes $D(r)$ for the
restriction with $r_i=b$ and all other coordinates free.
\begin{table}[ht]
\centering
\begin{tabular}{ccc}
\toprule
First queried coordinate $i$ & $D(x_i=0)$ & $D(x_i=1)$\\
\midrule
$0,4$ & $7/4$ & $11/4$\\
$1,2,3$ & $15/8$ & $21/8$\\
\bottomrule
\end{tabular}
\caption{Optimal remaining query costs under each first answer.
All other coordinates are free.}\label{tab:branches}
\end{table}
Both rows give
\[
 1+\tfrac12(7/4+11/4)
 =1+\tfrac12(15/8+21/8)=13/4.
\]
Here is a short derivation of the table. For $k\geq1$, the conjunction
$\operatorname{AND}_k$ of $k$ distinct fair bits, and likewise their
disjunction $\operatorname{OR}_k$, has optimal expected cost
$\sum_{j=0}^{k-1}2^{-j}$. Indeed, after any first query one answer determines
the value and the other leaves the same problem on $k-1$ bits. Thus the
costs for two and three bits are $3/2$ and $7/4$.

Two three-variable functions will also occur: $u(v\lor w)$ and
$u\lor vw$. Both have optimal cost $7/4$. Querying $u$ first leaves either
a constant or an OR of two bits (respectively an AND of two bits), so costs
$1+\tfrac12(3/2)=7/4$. Querying $v$ or $w$ first instead leaves a single
bit in one branch and an AND or OR of two bits in the other, for cost
$1+\tfrac12(1+3/2)=9/4$.

If $x_0=0$, the function is $x_1x_2x_3$, of cost $7/4$.
If $x_0=1$, write $(a,b,c,d)=(x_1,x_2,x_3,x_4)$; the function is
$abc\lor d(a\lor b\lor c)$. Querying $d$ first leaves AND$_3$ or OR$_3$,
giving cost $1+7/4=11/4$. Querying any of $a,b,c$ first leaves the two
three-variable types just considered, each of cost $7/4$, and also gives
$11/4$. This proves the first row.

If $x_1=0$, the remaining function has the form $uv(w\lor z)$. Querying
$u$ or $v$ first leaves a constant or a function of cost $7/4$, giving
$1+\tfrac12(7/4)=15/8$. Querying $w$ or $z$ first leaves AND$_3$ or
AND$_2$, giving $1+\tfrac12(7/4+3/2)=21/8$. Thus the minimum is $15/8$.
If $x_1=1$, the remaining function has the form $uv\lor wz$. Every first
query is equivalent by symmetry and leaves AND$_2$ in one branch and
$u\lor vw$ in the other. Its cost is therefore
$1+\tfrac12(3/2+7/4)=21/8$. This proves the second row.

Thus $D(*****)=a(f,\pi_{1/2})=13/4$. The full restriction table in
\texttt{data/certificate.json} also certifies these values, as described
in Section~\ref{sec:checks}.

\section{The partition cost}\label{sec:partition}
We first exhibit a partition cheaper than any algorithm, then give a
matching lower bound for all monochromatic partitions.
In Table~\ref{tab:partition}, each string lists coordinates in the order
$x_0,x_1,x_2,x_3,x_4$. The cells are pairwise disjoint and cover the cube,
and direct substitution in~\eqref{eq:f} gives the indicated constant value.
Here $|C|$ is the number of inputs in a cell $C$. Its total codimension
cost is $|C|\codim C$, the sum of that cell's cost over its inputs.
\begin{table}[ht]
\centering
\begin{tabular}{cccc}
\toprule
Cell $C$ & Value of $f$ & $|C|$ & $\codim C$\\
\midrule
\texttt{00***} & 0 & 8 & 2\\
\texttt{01*0*} & 0 & 4 & 3\\
\texttt{*111*} & 1 & 4 & 3\\
\texttt{0101*} & 0 & 2 & 4\\
\texttt{1*0*0} & 0 & 4 & 3\\
\texttt{110*1} & 1 & 2 & 4\\
\texttt{101*0} & 0 & 2 & 4\\
\texttt{10*11} & 1 & 2 & 4\\
\texttt{1*101} & 1 & 2 & 4\\
\texttt{10001} & 0 & 1 & 5\\
\texttt{11100} & 0 & 1 & 5\\
\bottomrule
\end{tabular}
\caption{A monochromatic partition with total codimension cost $102$.}
\label{tab:partition}
\end{table}
Consequently its expected codimension is
\begin{equation}\label{eq:upper}
 \frac{1}{32}\sum_C |C|\codim C
 =\frac{2\cdot8+3\cdot12+4\cdot10+5\cdot2}{32}
 =\frac{51}{16}.
\end{equation}
Together with the adaptive lower bound, this proves the strict inequality.

The partition also makes the distinction from a decision tree visible.
Every first query would have to fix its coordinate in every leaf cell.
Here $x_0$ is free in \texttt{*111*}, $x_1$ is free in \texttt{1*0*0},
and $x_2,x_3,x_4$ are free in \texttt{00***}. Thus no coordinate can be
the first query of a tree inducing this partition. This observation alone
does not prove the gap: another tree might still cost $51/16$,
and the adaptive lower bound rules that out.

For completeness, the partition is optimal. Encode a point by the integer
$k=\sum_{i=0}^4 2^i x_i$, and give it weight $y_k$ from
Table~\ref{tab:dual}. These weights sum to $102$ and satisfy
\begin{equation}\label{eq:dual}
 \sum_{x\in C} y_{\sum_i2^i x_i}\ \leq\ |C|\codim C
 \quad\text{for every monochromatic subcube }C.
\end{equation}
The inequalities are finite integer comparisons; both exact checkers
enumerate all restrictions and check~\eqref{eq:dual} on each monochromatic
one. No optimization solver is needed to verify them.
\begin{table}[t]
\centering
\begin{tabular}{c|rrrrrrrr}
\toprule
$k$ & 0 & 1 & 2 & 3 & 4 & 5 & 6 & 7\\
$y_k$ & $-1$ & 1 & 1 & 3 & 1 & 3 & 3 & 5\\
\midrule
$k$ & 8 & 9 & 10 & 11 & 12 & 13 & 14 & 15\\
$y_k$ & 1 & 3 & 3 & 5 & 3 & 5 & 5 & 3\\
\midrule
$k$ & 16 & 17 & 18 & 19 & 20 & 21 & 22 & 23\\
$y_k$ & 1 & 5 & 3 & 5 & 3 & 5 & 5 & 3\\
\midrule
$k$ & 24 & 25 & 26 & 27 & 28 & 29 & 30 & 31\\
$y_k$ & 3 & 5 & 5 & 3 & 5 & 3 & 3 & 1\\
\bottomrule
\end{tabular}
\caption{An integer dual certificate for the partition lower bound.}
\label{tab:dual}
\end{table}
Summing~\eqref{eq:dual} over any monochromatic partition $\mathcal C$
counts every input once, and gives
\[
 102=\sum_{k=0}^{31} y_k
 \leq\sum_{C\in\mathcal C}|C|\codim C.
\]
The negative weight at $k=0$ causes no difficulty: the cells form an exact
partition, not a cover with multiplicities. Division by $32$, together
with~\eqref{eq:upper}, proves $\scplx(f,\pi_{1/2})=51/16$ and completes
the proof of Theorem~\ref{thm:main}.

\section{Exact checks and formalization}\label{sec:checks}
This section describes the finite evidence for the displayed calculations
and distinguishes the Lean theorem from the additional partition lower
bound. The accompanying result archive contains the Lean project in
\texttt{lean/}, the certificate in \texttt{data/}, and the checkers in
\texttt{checks/}. File paths here are relative to that archive's result
directory; its \texttt{README.md} gives the build commands and a
fixed-version citation. We obtained the example using our own Python code
to enumerate monotone functions on five inputs up to coordinate
permutations. It computes adaptive costs by exact rational dynamic
programming and minimizes partition costs with Google OR-Tools' CP-SAT
solver~\cite{ORTools}. This discovery code is not included in the result
archive; the search and its optimizer are not needed to check the result.

The certificate in \texttt{data/certificate.json} contains the truth table,
a rational lower-bound table $L(r)$ for the $3^5=243$ restrictions, an
attaining decision tree, the partition, and the dual weights. Each of two
independently written Python checkers checks constancy and the inequalities
\[
 L(r)\leq0\quad\text{if $f$ is constant on $r$},\qquad
 L(r)\leq1+\frac{L(r^{i,0})+L(r^{i,1})}{2}
\]
for every free coordinate of every nonconstant restriction. Induction
using Lemma~\ref{lem:recurrence} gives $L(r)\leq D(r)$, in particular
$D(*****)\geq13/4$. Both checkers validate the decision tree and the
partition, check all the dual inequalities, and use exact rational
arithmetic. The second checker also recomputes the adaptive optimum
without relying on the lower-bound table. Corruption controls alter the
truth table, lower bounds, tree, partition, and dual weights; both checkers
must reject each altered certificate.

The Lean project proves the strict inequality independently of the Python
checkers and certificate. Its theorem \texttt{KS24.answer} states that
$f$ is monotone; every correct adaptive decision tree has total query cost
at least $104$ on the $32$ inputs; some such tree attains $104$; and some
monochromatic partition has total codimension $102$.
The theorem \texttt{KS24.gap} states that this partition costs strictly
less than every correct tree. Dividing the sums by $32$ gives precisely
$a(f,\pi_{1/2})=13/4$ and $\scplx(f,\pi_{1/2})\leq51/16$.
The lower bound $\scplx(f,\pi_{1/2})\geq51/16$ is established by the
separate dual certificate and the summation argument above; it is not
part of the Lean theorem.

The formal model counts queries only until the value is first determined.
It allows repeated queries, and proves its lower bound by induction on an
arbitrary finite decision tree. A deterministic algorithm truncated at
determination gives such a tree; conversely a tree can be extended by
querying all remaining bits without changing the cost up to determination.
A partition cell is represented by a bit mask recording which coordinates
are fixed, together with their values.
Every input must lie in exactly one cell, and codimension counts the fixed
coordinates among the five inputs. Discarding empty cells or mask bits
outside those coordinates leaves the same ordinary subcube partition.
These identifications connect the formal statement to the definitions
of~\cite{KS}.

Lean checks the finite calculations by reducing expressions inside its
proof-checking kernel, rather than trusting an external computation.
Build checks called axiom guards also fail if the final theorems acquire
assumptions beyond the listed logical axioms: \texttt{propext}
(propositional extensionality) and \texttt{Quot.sound} (soundness of
quotients). Sections~\ref{sec:alg} and~\ref{sec:partition} already give a
written proof of the strict inequality. Its independent Lean verification
depends on Lean's kernel and the correspondence of the finite definitions
with the mathematical ones. The exact partition
optimum additionally depends on the two Python checkers, Python's integer
and rational arithmetic, and the weak-duality argument above. The public
README gives the standard build and certificate-checking commands.

\begin{thebibliography}{9}
\bibitem{KS}
Laurin K\"ohler-Schindler and Jeffrey E. Steif,
\emph{A study of distributional complexity measures for Boolean functions},
ACM Transactions on Computation Theory (2026), published online May~19, 2026.
\url{https://doi.org/10.1145/3814613}.
First posted as arXiv:2408.12995 on August~23, 2024,
\url{https://arxiv.org/abs/2408.12995v1}.
Question and example numbering cited from arXiv version~1, the only version
listed in the arXiv history. The open author manuscript linked from Steif's
publications page retains Question~2.1 with the same wording and number:
\url{https://www.math.chalmers.se/~steif/p74.pdf}.
ACM's final typeset text was not accessible to us.

\bibitem{KRS}
Robin Kothari, David Racicot-Desloges, and Miklos Santha,
\emph{Separating Decision Tree Complexity from Subcube Partition Complexity},
Approximation, Randomization, and Combinatorial Optimization.
Algorithms and Techniques (APPROX/RANDOM 2015),
LIPIcs \textbf{40} (2015), 915--930.
\url{https://doi.org/10.4230/LIPIcs.APPROX-RANDOM.2015.915}.

\bibitem{ACKSW}
Srinivasan Arunachalam, Sourav Chakraborty, Michal Kouck\'y, Nitin Saurabh,
and Ronald de Wolf,
\emph{Improved bounds on Fourier entropy and Min-entropy},
arXiv:1809.09819, first posted September~26, 2018; version~2,
September~17, 2021.
\url{https://arxiv.org/abs/1809.09819}.
\bibitem{ORTools}
Google OR-Tools contributors,
\emph{OR-Tools: CP-SAT solver}.
\url{https://developers.google.com/optimization/cp/cp_solver}.

\end{thebibliography}
\end{document}
